\subsection{Soundness}\label{sec:soundness}

In \Cref{thm:soundness} we present the proof of soundness for tableaux as a calculus for satisfiability.
Our proof follows the same strategy used in~\cite{Fit83,Fit96,Priest08}. 
We start by introducing a well-known satisfiability preserving result regarding the use of new nominals.

% \begin{lem}
%     \label{lemma:nom}
%     \label{lemma:child}
%     Let $\varphi$ be a node expression and let $i$ be a nominal not in $\varphi$.
%     Then, $\varphi$ is satisfiable implies that $i{:}\varphi$ is satisfiable.
%     Moreover, if $\varphi = \psi \land \tup{@_i \pathh\alpha \cmpr \beta}$ and $j$ is a nominal not in $\varphi$, then, $\varphi$ is satisfiable implies $\psi \land (i{:}\tup{\pathh} j) \land (\tup{@_j \alpha \cmpr \beta})$ is satisfiable.
% \end{lem}

% \begin{proof}
%   Let $\model,n\vDash\varphi$. Since $i$ does not appear in $\varphi$, we can build a new model $\model'$ that is exactly like $\model$ with the exception that in $\model'$ we set $\naming(i)=n$. Thus, $\model',n\vDash i{:}\varphi$.
%   For the case $\varphi = \psi \land \tup{@_i \pathh\alpha \cmpr \beta}$, since $j$ is not in $\varphi$, we can use the same idea and set $\naming(j)=m$ for $m$ such that $(n,m)\in R_{\child}$.
% \end{proof}

% \fxnote{Escribo una nueva versión de esto porque no me resultaba claro. Fijense si es lo que hacía falta}

\begin{lem}
	\label{lemma:nom}
	\label{lemma:child}
        Let $\varphi$ be a node expression and let $j$ be a nominal not in $\varphi$.
		Then, (1) $\varphi$ is satisfiable iff $j{:}\varphi$ is satisfiable,
        (2) $\varphi = \psi \land i{:}\tup{\pathh}\chi$ is satisfiable iff $\varphi \land i{:}\tup{\pathh} j \land j{:}\chi$ is satisfiable, and 
        (3) $\varphi = \psi \land \tup{@_i \pathh\alpha \cmpr \beta}$ is satisfiable iff $\varphi \land i{:}\tup{\pathh} j \land (\tup{@_j \alpha \cmpr \beta})$ is satisfiable.
	% \end{enumerate}
\end{lem}

\begin{proof}
	In all cases, the implication from right to left is direct. 
	For (1), let $\model,n\vDash\varphi$. Since $j$ does not appear in $\varphi$, we can build a new model $\model'$ that is exactly like $\model$ with the exception that in $\model'$ we set $\naming(j)=n$. Thus, $\model',n\vDash j{:}\varphi$. 
	%The other direction follows since $j$ does not appear in $\varphi$. 
	%
	For (2) and (3), similarly, if we choose $j$ not in $\varphi$, we can use the same idea and set $\naming(j)=m$ for $m$ such that $(\naming(i),m) \in R_{\child}$. 
    %The other direction again follows since $j$ does not appear in $\varphi$.
\end{proof}


% The following lemma states that by decomposing a path expression using (child) preserves the satisfiability of a node expression.

% \begin{lem}
%     \label{lemma:child}
%     Let $\varphi = \psi \land \tup{@_i \pathh\alpha \cmpr \beta}$ be a node expression and $i$ a nominal not in $\varphi$.
%     Then, $\varphi$ is satisfiable implies $\varphi \land (i{:}\tup{\pathh} j) \land (\tup{@_j \alpha \cmpr \beta})$ is satisfiable.
% \end{lem}

% \begin{proof}
%     Similarly as in~\Cref{lemma:nom}, since $j$ does not appear in $\varphi$, we can set $\naming(j)=m$, for the corresponding $m$ such that $(n,m)\in R_{\child}$. 
% \end{proof}

Now we proceed to prove some properties relating the extension of a branch using tableaux rules, and the satisfiability of the node expressions contained in such a branch.

\begin{defi}
    \label{def:faithful}
    \label{def:branch-sat}
    Let $\Theta$ be branch of a tableau.
    We say that $\Theta$ is \emph{faithful} to a model $\model$ iff every node expression in $\Theta$ is satisfiable in $\model$.
    In turn, $\Theta$ is satisfiable iff there is a model $\model$ such that $\Theta$ is faithful to $\model$.
\end{defi}

The following lemma makes clear the way in which tableau rules work: they only introduce a contradiction (a clash) if such a contradiction is already present in the node expression we are building the tableau for.
Stated otherwise, tableau rules preserve satisfiability. 
%This is made precise in the following lemma.

\begin{lem}
    \label{lemma:soundness}
    The application of any one of the tableau rules in~\Cref{rules:basic,rules:nominals,rules:internalization,rules:path} to a satisfiable branch $\Theta$ yields another satisfiable branch.
\end{lem}

\begin{proof}
    Let $\Theta$ be a satisfiable branch of tableau,  i.e.,
    $\Theta$ is faithful to a model $\model$.
    The proof is by a case analysis of tableau rules.
    We develop some relevant cases.
    %The cases ($\land$) and $(\lnot\land)$ are immediate.
    For the rules ($\Diamond$) and (child), the result follows from \Cref{lemma:nom}.
    %
    For the rule (copy), notice that we have $\{i{:}\varphi, i{:}j\} \subseteq \Theta$ and that the extension after the application of (copy) is $\Theta' = \Theta \cup \{j{:}\varphi\}$.
    Moreover, we have $\model \vDash i{:}j$ and $\model \vDash i{:}\varphi$.
    From the first conjunct, we obtain $\naming(i) = \naming(j)$, and from the second $\model, \naming(i) \vDash \varphi$, 
    then $\model \vDash j{:}\varphi$.
    Therefore, $\Theta'$ is satisfiable.
    %
    %Let us now turn our attention to the case 
    For the rule (dRef),
    %Notice that in this case 
    we have $i{:}j \in \Theta$ and that the extension after the application of (dRef) is $\Theta' = \Theta \cup \{\tup{@_i =_{\cmp} @_j}\}$.
    Moreover, we have $\model \vDash i{:}j$, and thus $\naming(i) = \naming(j)$.
    Since $\approx_{\cmp}$ is an equivalence relation, $(\naming(i), \naming(j)) \in {\approx_{\cmp}}$, then $\model \vDash \tup{@_i =_{\cmp} @_j}$.
    Therefore, $\Theta'$ is satisfiable.
    %
    %The case (com$_1$) is also interesting. W.l.o.g. 
% %    Notice that in this case we have ${\cmpr} \in \{{=_{\cmp}}, {\neq_{\cmp}}\}$.
%  %   Let us 
%     consider $\cmpr$ is $=_{\cmp}$.
%     We have $\tup{@_i =_{\cmp} \beta} \in \Theta$ and that the extension after the application of (dRef) is $\Theta' = \Theta \cup \{\tup{\beta =_{\cmp} @_i}\}$.
%     Moreover, from the construction of the branch we have $\beta = @_j\beta'$, and from the assumption we have $\model \vDash \tup{@_i =_{\cmp} @_j\beta'}$.
%     The latter implies $\model \vDash \tup{@_i =_{\cmp} @_j\beta'}$.
%     From the definition of satisfiability, we know there is $\model, \naming(j), n \vDash \beta'$ s.t.\ $(\naming(i),n) \in {\approx_{\cmp}}$
%     Since $\approx_{\cmp}$ is an equivalence relation, $(n, \naming(i)) \in {\approx_{\cmp}}$.
%     From these results we get $\model \vDash \tup{\beta =_{\cmp} @_i}$.
%     This establishes that $\Theta'$ is satisfiable.
    Finally, notice that the cases (test), ($\lnot$test), ($\cup$), (${\lnot}{\cup}$), (int$_1$), and (int$_2$) are immediate from \Cref{prop:equivalences}.
\end{proof} 

%The lemma above states that our calculus preserves the truth of a node expression. 
The result in \Cref{lemma:soundness} enables us to state and prove the main result of this section.

\begin{thm}[Soundness]
    \label{thm:soundness}
    Let $\varphi$ be a node expression. If $\varphi$ is satisfiable, then there is a saturated branch of a tableau for $\varphi$ that is open.
\end{thm}

\begin{proof}
    %We prove the contrapositive, i.e., if all saturated tableaux for $\varphi$ are closed, then, $\varphi$ is not satisfiable. 
    Aiming for a contradiction, suppose all saturated tableaux for $\varphi$ are closed, but $\varphi$ is satisfiable. Any tableau for $\varphi$ starts with $i{:}\varphi$, for some $i\in\noms$ not appearing in~$\varphi$.
    From~\Cref{lemma:nom}, we know $i{:}{\varphi}$ is also satisfiable, thus the initial branch $\Theta=\set{i{:}\varphi}$ of the tableaux is satisfiable. 
    From~\Cref{lemma:soundness}, we know that the application of the rules in~\Cref{rules:basic,rules:nominals,rules:internalization,rules:path} always yields a satisfiable branch from $\Theta$.
    Eventually, we would obtain a saturated branch.
    This contradicts the assumption that all saturated tableaux built from $\varphi$ are closed.
    %Thus, $\varphi$ is not satisfiable.
\end{proof}



% \begin{thm}[Soundness]\label{thm:soundness}
%     For a finite set $\Sigma$ of node expressions, if $\Sigma \vdash \varphi$ then $\Sigma \vDash \varphi$.
% \end{thm}

% \begin{proof}
%     We prove the contrapositive. 
%     Supose that $\set{\varphi_1 \dots \varphi_n} \nvDash \varphi$. Then, there is a model $\model, m$ s.t.\ $\model, m \vDash \set{\varphi_1 \dots \varphi_n} \cup \set{\lnot \varphi}$. 
%     Then, $\chi = \varphi_1 \land \dots \land \varphi_n \land \lnot \varphi$ is satisfiable.
%     By \Cref{lemma:soundness}, any tableau for $\chi$ has a satisfiable branch $\Theta$.
%     We know that such a branch is not closed.
%     This implies there are no closed tableaux for $\chi$.
%     Finally, $\set{\varphi_1 \dots \varphi_n} \nvdash \varphi$.   
% \end{proof}